% Generated draft: requires mathematical and source review.
\documentclass[11pt]{article}
\usepackage[margin=25mm]{geometry}
\usepackage{fontspec}
\setmainfont{Latin Modern Roman}
\usepackage{amsmath,amssymb,mathtools,unicode-math}
\setmathfont{Latin Modern Math}
\usepackage{enumitem,needspace,xurl,hyperref}
\setlength{\parindent}{0pt}
\setlength{\parskip}{5pt}
\setlength{\emergencystretch}{3em}
\newcommand{\N}{\mathbb{N}}
\newcommand{\Z}{\mathbb{Z}}
\newcommand{\Q}{\mathbb{Q}}
\newcommand{\R}{\mathbb{R}}
\newcommand{\C}{\mathbb{C}}
\newcommand{\F}{\mathbb{F}}
\newcommand{\A}{\mathbb{A}}
\newcommand{\PP}{\mathbb{P}}
\newcommand{\E}{\mathbb{E}}
\newcommand{\eps}{\varepsilon}
\newcommand{\dd}{\,\mathrm{d}}
\begin{document}
\section*{UnsolvedMath ID: 20002050}
\textbf{Source problem code:} catalogue entry
\subsection{Definitions and mathematical statement}
\textbf{Problem 28.} When does the critical conformal operator $P=P_n$ on a closed connected even-dimensional Riemannian manifold satisfy $\ker P_n=\R$ on smooth functions? The source says this always holds for the two-dimensional Laplacian, while in dimension four there are both sufficient conditions and counterexamples, even among conformally flat metrics. Boundary conditions and connectedness must be explicit for the constants-only formulation.

\textbf{Supplied four-dimensional spectral criterion.} Use $\Delta=d^*d\ge0$ and the Paneitz convention
\[
 P_4=\Delta^2+\delta(T\,\dd),\qquad
 T=\tfrac23Rg-2\operatorname{Ric},\quad\delta=d^*.
\]
Let $\lambda_1>0$ be the first positive eigenvalue of $\Delta$ and let
\[
 \beta=\max\bigl(0,\max_{x\in M}[-\lambda_{\min}(T_x)]\bigr)
\]
be the least nonnegative constant with $T\ge-\beta g$. The aid asserts that if $\beta<\lambda_1$, then for every mean-zero smooth $u$,
\[
 \int_MuP_4u\,\dd V_g
 \ge\lambda_1(\lambda_1-\beta)\int_Mu^2\,\dd V_g.
\]
Thus $P_4$ is nonnegative and has exactly the constant kernel. The strict curvature-spectral margin is stated to be open in the $C^2$ topology on smooth metrics; this is a sufficient condition rather than a classification of all constant-kernel metrics.

For an Einstein four-metric the aid gives $P_4=\Delta(\Delta+R/6)$. Its nonconstant kernel is exactly the Laplace eigenspace of eigenvalue $-R/6$, when that number is a positive Laplace eigenvalue. For $R\ge0$ there is no such nonconstant resonance.

\textbf{Supplied conformally flat counterexample.} Let $k>0$ and take the product of the round sphere $S^2(k)$ of curvature $k$ with a closed connected hyperbolic surface $\Sigma(-k)$ of curvature $-k$. This four-metric is locally conformally flat and has scalar curvature zero. If $a$ and $b$ are Laplace eigenvalues of its two factors, separation of variables gives the Paneitz eigenvalue
\[
 (a+b)^2-2ka+2kb.
\]
The aid asserts that the kernel consists of constants and the three first spherical harmonics pulled back from $S^2(k)$, hence has dimension four. Here $\beta=2k$ and $\lambda_1(M)=\min(2k,\lambda_1(\Sigma))$. Whenever $\lambda_1(\Sigma)\ge2k$, this gives $\beta=\lambda_1(M)$ while the kernel is nonconstant, so replacing the strict sufficient inequality by a non-strict one is invalid. The summary associates this example with Singer--Eastwood but supplies no full reference.

\textbf{Adjacent determinant material, not part of the kernel question.} Section III is extracted from T.~Branson and R.~Gover's ``Origins, applications, and generalizations of the $Q$-curvature''. It discusses a natural operator $A$ with positive definite leading symbol, a positive power of a conformally invariant operator, such as a GJMS operator or a Dirac square. In dimensions two, four and six, it reports the schematic anomaly formula
\[
 -\log\frac{\det_\zeta\widehat A}{\det_\zeta A}
 =\alpha\left\{\frac12\int_M\omega P\omega\,\dd V_g
                  +\int_M\omega Q\,\dd V_g\right\}
       +\mathcal C_F[\widehat g,g]+H.
\]
Here $\widehat g=e^{2\omega}g$, $\mathcal C_F$ is a difference of integrated local scalar invariants, and $H$ is a nullspace-dependent correction, stated to vanish when $\ker A=0$. The hats/order in the extracted $F$ terms are damaged, so their sign is not guessed. Zeta determinant conventions with zero modes or negative eigenvalues require specification; positivity of the principal symbol alone does not supply positivity or invertibility of $A$. In particular $\ker P_4=\R$ is not the same condition as $\ker A=0$.
\subsection{Short English statement}
The source asks when a critical conformal operator has only constant solutions. The supplied aid gives a strict curvature-versus-Laplace spectral bound for Paneitz, an Einstein resonance criterion and a conformally flat four-dimensional counterexample. The appended determinant discussion is separate and uses a different zero-nullspace condition.
\subsection{Sources}
\begin{itemize}
\item Local UnsolvedMath registry, permanent ID 20002050, source code catalogue entry; canonical \texttt{aim-geometry-notes.json}, notes[387].
\item AIM \emph{Conformal structure in geometry, analysis, and physics}, Problem 28 and beginning of Section III: \url{https://aimath.org/WWN/confstruct/confstruct.pdf}.
\item T.~Branson and R.~Gover, ``Origins, applications, and generalizations of the $Q$-curvature'', named as the source of the adjacent determinant discussion, equation (40). No full publication data are supplied. Singer--Eastwood attribution appears in the aid for the product example, without a citation.
\item Registry research attempt 1, dated 2026-08-14: partial result, rigor ``proved'', novelty confidence low. The machine-generated, non-peer-reviewed aid requires independent verification before citation.
\end{itemize}
\end{document}
